% Generated draft: requires mathematical and source review.
\documentclass[11pt]{article}
\usepackage[margin=25mm]{geometry}
\usepackage{fontspec}
\setmainfont{Latin Modern Roman}
\usepackage{amsmath,amssymb,mathtools,unicode-math}
\setmathfont{Latin Modern Math}
\usepackage{enumitem,needspace,xurl,hyperref}
\setlength{\parindent}{0pt}
\setlength{\parskip}{5pt}
\setlength{\emergencystretch}{3em}
\newcommand{\N}{\mathbb{N}}
\newcommand{\Z}{\mathbb{Z}}
\newcommand{\Q}{\mathbb{Q}}
\newcommand{\R}{\mathbb{R}}
\newcommand{\C}{\mathbb{C}}
\newcommand{\F}{\mathbb{F}}
\newcommand{\A}{\mathbb{A}}
\newcommand{\PP}{\mathbb{P}}
\newcommand{\E}{\mathbb{E}}
\newcommand{\eps}{\varepsilon}
\newcommand{\dd}{\,\mathrm{d}}
\begin{document}
\section*{UnsolvedMath ID: 6200081}
\textbf{Source problem code:} catalogue entry
\subsection{Definitions and mathematical statement}
Let $G\leq\operatorname{Isom}(\mathbb H^n)$ be discrete, torsion-free, and finitely generated, with no abelian subgroup of rank at least two. Let $\Lambda(G)\subset\partial_\infty\mathbb H^n$ be the orbit accumulation set. Its conical subset $\Lambda_c(G)$ consists of endpoints $\xi$ for which an orbit has infinitely many points within a fixed bounded distance of a geodesic ray toward $\xi$, tending to $\xi$. Compute Hausdorff dimension using the standard round boundary metric, and put $\delta_c=\dim_H\Lambda_c(G)$. Let $\operatorname{cd}_{\mathbb Z}(G)$ be integral cohomological dimension, the projective dimension of the trivial $\mathbb ZG$-module $\mathbb Z$.

The source asks:
\begin{enumerate}
\item Is $\operatorname{cd}_{\mathbb Z}(G)\leq\delta_c+1$? It reports an affirmative answer with homological dimension in place of cohomological dimension, citing its reference [36].
\item If equality holds, must $\Lambda(G)$ be a round sphere? The supplied text continues ``and $G$?'' without its predicate, so any further conclusion about $G$ is missing. It reports the intended rigidity statement known for geometrically finite $G$, again citing [36].
\item If $\delta_c<2$, must $G$ be geometrically finite?
\item If $\delta_c<1$, must $G$ be a classical Schottky-type group?
\end{enumerate}
Here geometrical finiteness may be expressed by finite volume of a fixed-radius neighborhood of the convex core. A round sphere in the boundary may be a round subsphere; the record does not state its dimension. The classical Schottky model uses generators pairing disjoint round boundary balls, but the source's ``Schottky-type'' terminology is undefined, especially concerning parabolics or tangencies. That qualifier is not removed. Elementary groups and the dimension convention for an empty conical limit set are also unspecified.
\subsection{Short English statement}
Relate integral cohomological dimension and geometric finiteness of a torsion-free Kleinian group to the Hausdorff dimension of its conical limit set. Retain the equality-rigidity and Schottky-type questions, while flagging a truncated conclusion and missing definitions.
\subsection{Sources}
\begin{itemize}
\item Misha Kapovich, \emph{Problems on Boundaries of Groups and Kleinian Groups} (2005), Problem 81, PDF page 20, \url{https://www.math.ucdavis.edu/~kapovich/EPR/problems.pdf}; accessed 2026-07-29, PDF-text-normalized extraction.
\item Source reference [36], cited for the homological-dimension inequality and geometrically finite equality case, is not identified in the supplied record. The research summary also mentions Bestvina--Mess and a Bishop--Jones-type result without full citations or exact hypotheses.
\item Local registry \texttt{data/unsolvedmath/problems.json}, permanent UnsolvedMath ID $6200081$, source code catalogue entry, \emph{Boundaries of Groups and Kleinian Groups --- Problem 81}. Dataset attribution: Ulam AI and the UnsolvedMath Contributors.
\end{itemize}
\end{document}
